\cvsection{%
  \awesomeFontNew \symbol{"F275}
  \awesomeFontNew \symbol{"F2B5}
  \awesomeFontNew \symbol{"F0D6}%
}{Work Experience}

\vspace{-2mm}

\begin{xltabular}{\textwidth}{r@{}X@{}}

% MICROGRAPHIA BIO -----------------------------------------------
\multicolumn{1}{@{}r@{}|}{%
  \begin{tabular}{r} \emph{Curr.} \\ \\ \\ \\ \\ 2025 \end{tabular}%
}&%
\multicolumn{1}{l}{%
  \begin{tabularx}{0.9\textwidth}{@{}X@{}}
    {\headingfont\textbf{Staff Engineer, Platform \& Systems}},
    \small at \href{https://www.micrographiabio.com}{Micrographia Bio}, London, UK.
    \hfill \small{\color{darkgrey}Python · Pytorch . Ansible · AWS }
    \vspace{2mm} \\
    \small\textit{AI-first life sciences company building a proprietary, end-to-end spatial proteomics platform for drug discovery, developed entirely in-house.}
    \vspace{1.5mm} \\
    \small
      \tabitem Deliver cross-functional features across the full platform stack, tracing how bench changes propagate through the imaging pipeline to spatial proteomics results.
      \tabitem Coach-player across a multidisciplinary team (biochemistry, optics, robotics, software): pairing, code review, and 1:1s that keep research and engineering aligned.
      \tabitem Own observability and DevSecOps: experiment QA/QC, environment standardisation, containerisation, and infrastructure automation.
      \tabitem Own quarterly EC2 spend audits, deprecating, deleting, or archiving to Glacier for low-hanging savings; cut AWS costs and keep cloud spend intentional and defensible.%
  \end{tabularx}%
}\\
\multicolumn{2}{c}{}\vspace{1mm}\\

\multicolumn{1}{@{}r@{}|}{%
  \begin{tabular}{r} 2025 \\ \\ \\ \\ 2024 \end{tabular}%
}&%
\multicolumn{1}{l}{%
  \begin{tabularx}{0.9\textwidth}{@{}X@{}}
    {\headingfont\textbf{Platform Engineer}},
    \small at \href{https://www.micrographiabio.com}{Micrographia Bio}, London, UK.
    \hfill
    \vspace{2mm} \\
    \small
      \tabitem Product owner and technical lead for the company's core IP in spatial proteomics, built with no prior art; coached a multidisciplinary team (optics, chemistry, biology, robotics, software) through pairing, code review, and 1:1s to align on architecture decisions.
      \tabitem Productionised R\&D through end-to-end experiment design, standardising workflows from cell preparation through imaging, ML analysis, and result interpretation.%
  \end{tabularx}%
}\\
\multicolumn{2}{c}{}\vspace{1mm}\\

\multicolumn{1}{@{}r@{}|}{%
  \begin{tabular}{r} 2024 \\ \\ \\ \\ \\ 2021 \end{tabular}%
}&%
\multicolumn{1}{l}{%
  \begin{tabularx}{0.9\textwidth}{@{}X@{}}
    {\headingfont\textbf{Computer Vision Research Engineer}},
    \small at \href{https://www.micrographiabio.com}{Micrographia Bio}, London, UK.
    \hfill 
    \vspace{2mm} \\
    \small
      \tabitem Joined as an early CV/ML hire and grew into a senior member of the software team; early contribution to productionisation, building MLOps for automated microscopy and scientific image processing and cutting day-long experiments to under an hour (incl.\ a 141\,$\to$\,19 min pipeline step, 7$\times$ speedup).
      \tabitem Worked in close collaboration with the research team, evolving the existing infrastructure into a plugin-based framework that gave R\&D consistent, modular tooling for task and experiment design, enabling rapid iteration across the full pipeline.
      \tabitem Applied that framework to adapt transformer and heavy vision models for our spatial proteomics drug discovery use cases, running 90+ evaluations in two months to identify the best-fit architectures.%
  \end{tabularx}%
}\\
\multicolumn{2}{c}{}\vspace{1mm}\\

% VOLUNTEER WORK -------------------------------------------------
\cventry{2020}{2021}
  {Open-source Volunteer}
  {Anonymous Contribution, Remote}
  {Python · Edge ML}
  {%
    \tabitem Architected CV-based social distancing monitoring and MLOps pipeline, deployed to local businesses and individuals worldwide to help navigate the pandemic.
    \tabitem Optimised, tested, and released models for Edge ML deployment on embedded hardware (Jetson, Raspberry Pi, Coral TPU).%
  }

% ARSAM ROBOTICS -------------------------------------------------
\cventry{2020}{2021}
  {AI Team Lead}
  {\href{http://www.arsamrobotics.com/en/index}{\textsc{Arsam Robotics}}, Remote}
  {Python · Keras · TFLite · Android · OpenCV}
  {\small\textit{Robokids: interactive robotic amusement park and kids club.}\\[1mm]
    \tabitem Co-founded a ``phygital'' educational board game (\href{https://www.unicef.org/iran/en/press-releases/unicef-and-vice-presidency-science-and-technology-hold-national-innovations-award}{\textbf{UNICEF Innovation Award}}): real-time CV reads physical tokens on a custom mat, driving a parent-customisable digital learning world on mobile.
    \tabitem Led and grew the AI team from concept to product; designed CV architecture, overcame zero training data via synthetic bootstrapping and semi-automated annotation; optimised models for real-time mobile inference (TFLite/Android).%
  }

% TEHRAN INTERNET ------------------------------------------------
\cventry{2019}{2021}
  {Head of AI Department}
  {\href{https://www.linkedin.com/company/tehran-internet}{\textsc{Tehran Internet Holdings}}, Remote}
  {Python · Keras · TFLite · Java · Android}
  {\small\textit{Large Iranian internet and e-commerce holding company, operating across fintech, retail, and logistics.}\\[1mm]
    \tabitem Founded the organisation's first AI team; grew it from CV applications into customer segmentation, clustering, and recommender systems.
    \tabitem Built a high-performance mobile library for real-time bilingual payment card scanning (\href{https://780.ir/}{*780\#}), handling non-standard multi-bank card designs via synthetic bootstrapping and semi-automated annotation.
    \tabitem Built a SIM card scanning app for the \href{https://www.eways.co/}{\textsc{Eways}} storehouse, achieving 99.98\% accuracy in the first year; still in production today.%
  }

% BAYAN ----------------------------------------------------------
\cventry{2018}{2019}
  {Computer Vision Engineer}
  {\href{https://vision.bayan.ir}{\textsc{Bayan}}, Tehran, Iran}
  {Python · Qt · TensorFlow · PyTorch · Caffe · MXNet}
  {\small\textit{Information technology and consulting company, Tehran.}\\[1mm]
    \tabitem Built a real-time pedestrian detection system with a synthetic dataset (Google Street View + Mask-RCNN) to cover Iranian street clothing; owned the multi-framework codebase and Linux infrastructure.%
  }

% FARAADID -------------------------------------------------------
\cventry{2016}{2017}
  {Computer Vision Engineer \& Head of DevOps}
  {\href{https://faraadid.com/}{\textsc{Faraadid}}, Tehran, Iran\quad\href{https://github.com/Faraadid/}{\seticon{faGithub}}}
  {C++ · Python · Bash · Qt · Jenkins · GitLab}
  {\small\textit{Early-stage Computer Vision startup, Tehran.}\\[1mm]
    \tabitem Architected a network-configurable CV Engine in C++, developed internally, that became the heart of 20+ industrial projects delivered by the company.
    \tabitem Built DevOps and internal engineer-experience tooling from scratch to accelerate R\&D.
    \tabitem Core team member on \textbf{FaraTraffic}, the company's most successful product: contributed to system architecture, deployment, and load balancing, alongside computer vision and algorithm development for vehicle detection and traffic analysis; 5+ years of production uptime.%
  }

% EARLIER ROLES --------------------------------------------------
\cventry{2011}{2016}
  {Software Engineer (various roles)}
  {Tehran \& Isfahan, Iran}
  {Java · C++ · Linux · Android · Assembly}
  {%
    \tabitem \textsc{Khak} (2014): engineered the engine of \href{http://khatkhan.com/}{\textsc{Khatkhan}}, a popular Persian e-book reader.
    \tabitem \textsc{CVAS} (2013): architected an e-book engine deployed in \textsc{Ketabrah}, \textsc{Behketab}, and \textsc{Toranj}.
    \tabitem Robotics Lab, Isfahan (2012): built a face vectorisation system rendered by a 6-DOF robot arm (Mitsubishi RV1-A).
    \tabitem \href{https://www.jooyeshgar.net/}{\textsc{Jooyeshgar}} (2011): developed embedded Linux systems: DVR and 3D printer.%
  }

\end{xltabular}
